% r06 The heavy shell II: the heavy-shell estimates and their continuum form
\section{The heavy shell II: the heavy-shell estimates and their continuum form}

\subsection{The heavy-shell estimates}

In the curvature sector \REz\ is proved at one configuration, from two estimates on King's walled torus model with the shell mass inside the matter covariance, the heavy-shell estimates \HS\ of \Cref{rot:hyp:HS}. The shell is now taken to be the whole exterior of the source box in the torus. Let $(\fs,\chi)$ be a two-level configuration in the sense of \Cref{rot:def:shell,rot:def:twolevel}, with the walled torus $T\defeq D_{J'}$ of $\fc=(J',\psi,n)$ and the source box $\Hboxp\defeq\Hbox=\Qb{b'}$, in which the shell $C_L\setminus\Hbox$ is replaced by
\[
U\defeq T\setminus\Hboxp ,
\]
so that $\partial U=\partial\Hboxp$ is a single convex surface with no re-entrant edge (on a cubic shell $C_L\setminus\Hbox$ the derivative clause of King's bound on the block propagators, \cite[Thm.~3.3]{K1}, fails at the re-entrant edges); the plateau condition $\psi\ge1$ on $\Nr1U$ is kept, the profile is the sharp one, $\chi=\mathbf 1_U$, which \Cref{rot:rem:Eproof}(ii) allows for every lattice statement of the preceding section, so that the dead set is $\Dchi=U$, and $\tauGB(\fs)$ is computed, as in \Cref{rot:def:shell}, from $\Hbox$ and a King cube $C_L$ with $\Hboxp\subset C_L\subset T$ at $\ell^\infty$-distance $D(\fs)+3$ from $\Hboxp$. For $\varkappa\in[0,\infty)$ let $\Psik(n,\varkappa)$ be the two-level expectation of \Cref{rot:def:twolevel} with the shell mass $\varkappa\mathbf 1_U$ added to the wall: by \Cref{rot:def:tori}(c) the wall $n\psi+\varkappa\mathbf 1_U$ raises the scalar mass squared by $2n\psi$ on the plateau and by a further $2\varkappa$ on $U$, and $\Psik(n,K)$, $\Psik(n,\infty)$ of \Cref{rot:hyp:RES} are the cases $\varkappa=K$ and $\varkappa=\infty$ (the matter killed on $U$). Let $Z_\varkappa(zf,0)$ be the corresponding partition function with the curvature source $zf$, so that $Z_\varkappa(zf,0)/Z_\varkappa(0,0)=\langle e^{zF_{\Lk}(f)}\rangle_\varkappa$ is the expectation of the source weight in the walled torus state with the shell mass $\varkappa\mathbf 1_U$. Here $Z_\varkappa(\,\cdot\,,0)$ is the unnormalised functional of this walled state in the convention of \Cref{rot:def:Zun}, with the wall $n\psi$ and the shell mass $\varkappa\mathbf 1_U$ in its single-site measures; it is not the normalised functional $Z^a$ of the first sections, and \Cref{rot:thm:shell} uses it only through the ratios $Z_\varkappa(zf,0)/Z_\varkappa(0,0)$. Finally $n_*\defeq1$: the estimates of \HS\ are read at plateau heights $n\ge n_*$, at which the mass squared on the plateau is raised by at least $2$.

The first statement carries the three steps of the dead-shell chain, proved in the preceding section for the cubic shell $C_L\setminus\Hbox$, to the exterior-shell geometry.

\begin{proposition}[\DSh: the dead-shell chain for the exterior shell]\label{rot:hyp:DS}
For the two-level configurations with the exterior shell $U=T\setminus\Hboxp$ the three steps of the proof of \Cref{rot:thm:SHtoTIGHT} hold: pinning the photon outside a King cube $C_L\supset\Hboxp$ all of whose sites outside $\Hboxp$ are dead, the photon switch across the matter-free collar between $\Hboxp$ and $\partial C_L$, and the passage to the comparison cube on $X_L=C_L\cap\Lk$. Consequently the dead-shell bound \eqref{rot:eq:deadbound} holds for them: at every $k$ with $\eps_k$ small, $\Psik(n,\infty)\le Z_{k,L}(\mathrm if)+(e^{\tauGB(\fs)}-1)^{1/2}$ at $g=0$.
\end{proposition}

% r06_fill_DS: proof of Proposition rot:hyp:DS, the dead-shell chain in the exterior-shell geometry.
% Input from sec/r06_rot_shell2.tex right after the statement of \Cref{rot:hyp:DS}.

\begin{proof}[Proof of \Cref{rot:hyp:DS}]
Fix a two-level configuration with the exterior shell: the walled torus $T=D_{J'}$ of $\fc=(J',\psi,n)$, the source box $\Hboxp=\Qb{b'}$ with $b'\ge b+3$, the sharp profile $\chi=\mathbf 1_U$ with $U=T\setminus\Hboxp$, so that the dead set is $\Dchi=U$, and a King cube $C_L$ with $\Hboxp\subset C_L\subset T$, $\dist_\infty(\Hboxp,\R^3\setminus C_L)=D(\fs)+3$, $D(\fs)\ge\max(1,3/\lmu)$ and $\Nr1{\bar C_L}\subset(-25MJ',25MJ')^3$; let $f$ be a real $2$-form supported in $\Qb b$, $g=0$, and $k$ with $\eps=\eps_k\le\min(1/\lmu,1/4)$. Write $\Theta\defeq e^{\mathrm iF_{\Lk}(f)}$ and $E\defeq T\setminus C_L$ for the set of sites of the torus outside $C_L$. Since $\supp f\subset\Qb b$ and $b'\ge b+3$, every plaquette of $f$ lies in $\Hboxp$ at $\ell^\infty$-distance $>1$ from $\partial\Hboxp$. The dead-shell state $\Psik(n,\infty)=\langle\Theta\rangle$ is the walled torus state of $\fc$ with the matter killed on $U\cap\Lk$ (\Cref{rot:def:twolevel}): the single-site measure at every site of $U$ is the point mass at $\phi=0$, every hopping term with an endpoint in $U$ is reduced to its diagonal part at the other endpoint, and the photon is the periodic photon of $T$ with mass $\lmu$ (\Cref{rot:def:tori}(d)). The matter alive is that of $\Hboxp\cap\Lk$, with the wall $n\psi$ in its single-site measures; since $\psi=0$ on $\Nr1{\Qb b}$ and $\psi\ge1$ on $\Nr1U$, the wall is a fixed non-negative function on $\Hboxp$. The two differences from the cubic shell of \Cref{rot:def:shell} are that the dead region $U$ now extends to the whole exterior of $\Hboxp$ in $T$, so that no live matter sits outside $\Hboxp$, and that $\partial U=\partial\Hboxp$ is a single convex surface. Both only simplify the three steps.

\emph{Step 1 (pinning).} For $s\ge0$ add $s|A_\beta(y)|^2$ to the action for every photon variable $A_\beta(y)$ on a bond with an endpoint in $E$, that is for every $(\beta,y)$ with $y\in E$ or $y+\eps e_\beta\in E$. The resulting measures are BFS measures of King type (the paragraph after \Cref{rot:def:Zun}: adding $s|\ell(A)|^2$ to the action is one of the admitted operations), so by \Cref{lem:s02:BFS}(a) the expectation of $\Theta$, which is of class $\cI$, does not decrease in $s$; let $s\to\infty$. In the limit the pinned variables are fixed at $0$, and by the Gaussian structure of the photon part the law of the remaining variables $A_\alpha(x)$, $x,x+\eps e_\alpha\in C_L$, is the periodic photon of $T$ conditioned on the vanishing of all the others: this is the law ($\natural$) of \Cref{rot:lem:GB}, whose cell contains $\Nr1{C_L}$ by hypothesis. In the pinned measure the matter outside $\Hboxp$ is dead, so the hopping terms with both endpoints in $U$ contribute nothing (their cosine parts vanish at $\phi=0$ and their diagonal parts vanish as well), and a hopping term between a site of $\Hboxp$ and a dead neighbour in $U$ is its diagonal part $\frac\eps2r_x^2$ at the live endpoint, a Dirichlet term independent of $A$. In particular no hopping term couples $E$ to a live site, since every site of $C_L$ adjacent to $E$ lies in $C_L\setminus\Hboxp\subset U$ and is dead, and the matter on $E$ is a product of point masses, a factor $1$. What remains is exactly the $\natural$-measure of \Cref{rot:lem:GB}: the matter on $\Hboxp\cap\Lk$ with BFS single-site measures (the wall $n\psi$ and the Dirichlet terms toward $U$ included), hopping terms only between neighbouring sites of $\Hboxp$, coupled to the photon variables of the bonds of $\Hboxp$ as in King's action, and the photon law ($\natural$) on the bonds of $C_L$. Hence
\[
\Psik(n,\infty)\ \le\ \langle\Theta\rangle_\natural .
\]

\emph{Step 2 (the photon switch across the collar).} The collar $C_L\setminus\Hboxp\subset U$ carries no live matter, and $\dist_\infty(\Hboxp,\R^3\setminus C_L)=D(\fs)+3$ with $D(\fs)\ge\max(1,3/\lmu)$ and $\eps\le\min(1/\lmu,1/4)$. So \Cref{rot:lem:GB} applies to the pair of measures with the common matter part of Step~1 and the photon parts ($\natural$) and (N), the latter King's Neumann photon on $C_L\cap\Lk$; its hypotheses on the observable hold with $c=0$ and $g=0$, $f$ being supported on plaquettes of $\Hboxp$. It gives
\[
\langle\Theta\rangle_\natural\ \le\ \langle\Theta\rangle_N+\langle|\Theta|^2\rangle_N^{1/2}\big(e^{\tau}-1\big)^{1/2}
=\langle\Theta\rangle_N+\big(e^{\tau}-1\big)^{1/2},\qquad
\tau\le C_{\rm GB}\,\srw\,\vol_\eps(\Nr2{\Hboxp})\,D(\fs)^{-1}e^{-\frac32m_\eps D(\fs)},
\]
since $|\Theta|^2=1$ at $g=0$. With $m_\eps\ge0.96\lmu$ and $\vol_\eps(\Nr2{\Hboxp})\le\vol(\Nr3{\Hboxp})$ for $\eps\le1$ this is $\tau\le\tauGB(\fs)$, the constant \eqref{rot:eq:tau} computed from $\Hbox=\Hboxp$ and $D(\fs)$, as in Step~2 of the proof of \Cref{rot:thm:SHtoTIGHT}. The geometry of the shell outside $C_L$ plays no role here: the lemma sees only the matter on $\Hboxp$ and the photon on $C_L$.

\emph{Step 3 (to the comparison cube).} The $N$-measure of Step~2 is the Neumann--Dirichlet state on $X_L=C_L\cap\Lk$ of \Cref{rot:def:tori}(a) with the wall $n\psi$ in its single-site measures and the matter killed on $(C_L\setminus\Hboxp)\cap\Lk$: King's Neumann photon on $C_L$, Dirichlet terms for the matter of $\Hboxp$ toward its dead neighbours and, for the sites of $C_L$ adjacent to $\R^3\setminus C_L$, toward the outside (these sites are dead, so the two readings of their Dirichlet terms coincide). By \Cref{lem:s02:Kprime}, applied with the set $(C_L\setminus\Hboxp)\cap\Lk$ and the rate $1$ on it, this state is the limit $s\to\infty$ of the walled Neumann--Dirichlet states on $X_L$ in which the single-site measures on $(C_L\setminus\Hboxp)\cap\Lk$ are multiplied by $e^{-s\eps^3r_x^2}$; $|\Theta|\le1$ is of the form admitted there. By \Cref{lem:s02:BFS}(c) the expectation of $\Theta$ does not decrease when $s$ decreases to $0$ (the single-site weights $e^{-s\eps^3r_x^2}$ are removed), nor when the wall factors $e^{-\eps^3n\psi(x)r_x^2}$, $x\in X_L$, are then removed (the same clause, the wall being a non-negative multiplier of the form $e^{-sr^2}$ at each site). The result is the comparison cube $\ava\cdot$ on $X_L$ (\Cref{def:s02:acube}): the Neumann photon on $C_L$, the scalar alive on all of $X_L$ with the Dirichlet condition outside $C_L$, no wall. In it $\Theta$ is the same function of the shared variables, every plaquette of $f$ lying in $\Hboxp$ at $\ell^\infty$-distance $>1$ from $\R^3\setminus C_L$. Hence $\langle\Theta\rangle_N\le\ava{e^{\mathrm iF(f)}}|_{J=L}=Z_{k,L}(\mathrm if)$, and with Steps~1--2
\[
\Psik(n,\infty)\ \le\ Z_{k,L}(\mathrm if)+\big(e^{\tauGB(\fs)}-1\big)^{1/2},
\]
which is \eqref{rot:eq:deadbound} at $g=0$ for the exterior-shell configuration. The three steps use the shape of $U$ only through two facts, that every site of $C_L$ adjacent to $E$ is dead (Step~1) and that the collar between $\Hboxp$ and $\partial C_L$ carries no live matter (Step~2); both hold because $U\supset C_L\setminus\Hboxp$. The two-level sandwich \Cref{rot:lem:TL}, which \Cref{rot:thm:REStoTIGHT} applies to this configuration, needs $\chi\le\psi$ pointwise, that is $\psi\ge1$ on $U=T\setminus\Hboxp$; this is the plateau condition $\psi\ge1$ on $\Nr1U$ of \Cref{rot:def:shell}, kept for the exterior shell, and it is compatible with $\psi=0$ on $\Nr1{\Qb b}$ because $\dist_\infty(U,\Qb b)\ge3$.
\end{proof}

With \Cref{rot:hyp:DS} in place, the proof of \Cref{rot:thm:REStoTIGHT} runs verbatim for the exterior-shell configurations of \Cref{rot:thm:shell}: its only use of the shell geometry is the dead-shell bound \eqref{rot:eq:deadbound}, which the proposition supplies, and \Cref{rot:lem:TL}, which holds for the sharp profile by \Cref{rot:rem:Eproof}(ii).


\begin{theorem}[\HS: the heavy-shell estimates]\label{rot:hyp:HS}
Let $(\fs,\chi)$ be a two-level configuration with the exterior shell $U=T\setminus\Hboxp$, with $\Hboxp$ and $T$ unions of blocks of the final scale of King's block-spin flow, with $\Hboxp$ of side at least two final blocks and at distance at least three final blocks from its periodic images in every direction, and with every corner coordinate of $\Hboxp$, counted from the origin in units of the final block, congruent to $M'/2$ modulo $M'$, where the large-block integer is $M=5M'$ with $M'$ even (the grid-offset condition; King's tori satisfy the clauses on $T$, and in each direction every interval of $M'$ final blocks contains an admissible corner coordinate) and plateau height $n\ge n_*$, let $g=0$ and let $f$ be a real $2$-form in $C^\infty_c$ supported in $\Qb b$. There are $\eta_1>0$ and, for every stopping scale $\eta_0\in5^{-\N}$ with $\eta_0\le\eta_1$: numbers $\sfa(\eta_0),\sfb(\eta_0)\in[0,\infty)$ with $\sfa(\eta_0),\sfb(\eta_0)\to0$ as $\eta_0\to0$; a threshold $\varkappa_0(\eta_0)<\infty$; a spacing bound $\eps_0(\eta_0)>0$; a function $r(\,\cdot\,;\eta_0)$ on $[\varkappa_0,\infty)$ with $r(\varkappa;\eta_0)\to0$ as $\varkappa\to\infty$; and, for every $k$ with $\eps_k\le\eps_0(\eta_0)$, a measure space $(\tilde\chi,db)$ not depending on $\varkappa$ and measurable functions $\tilde S_\varkappa(\,\cdot\,;z)$ on $\tilde\chi$ for $\varkappa\ge\varkappa_0$ and $z\in[-1,1]$ (both may depend on $k$; the index is suppressed); such that, with
\[
\Icap_\varkappa(z)\defeq\int_{\tilde\chi}e^{-\tilde S_\varkappa(b;z)}\,db ,\qquad
\Delta_z(b)\defeq\tilde S_{\varkappa'}(b;z)-\tilde S_\varkappa(b;z),
\]
for every such $k$, all $\varkappa_0\le\varkappa\le\varkappa'<\infty$ and all $z\in[-1,1]$:
\begin{enumerate}[label=\textup{(\roman*)},leftmargin=*]
\item \textup{(the two-sided bound)} $e^{-\sfa}\,\Icap_\varkappa(z)\ \le\ Z_\varkappa(zf,0)\ \le\ e^{\sfa}\,\Icap_\varkappa(z)+\sfb\,Z_\varkappa(0,0)$;
\item \textup{(the $\varkappa$-Cauchy property)} $\sup_{b,b'\in\tilde\chi}\,|\Delta_z(b)-\Delta_0(b')|\ \le\ 2\,r(\varkappa;\eta_0)$.
\end{enumerate}
\end{theorem}

% r06_fill_HS: the heavy-shell estimates (HS), Theorem rot:hyp:HS: the two shell theorems and the proof of (HS) from them.
% Input from sec/r06_rot_shell2.tex in place of the proof of \Cref{rot:hyp:HS}.
\providecommand{\Ish}{\ensuremath{I^{\sharp}}}
\providecommand{\Ifull}{\ensuremath{I}}
\providecommand{\Scap}{\ensuremath{\tilde S^{\mathrm{cap}}}}
\providecommand{\chicap}{\ensuremath{\tilde\chi^{\mathrm{cap}}}}
\providecommand{\kz}{\ensuremath{\varkappa_0}}

\subsubsection*{The shell model and the stopping-scale integrals}

Fix the exterior-shell configuration of \Cref{rot:hyp:HS}: the torus $T=D_{J'}$, the box $\Hboxp=\Qb{b'}$, both aligned with the blocks of the final scale as \Cref{rot:hyp:HS} requires (block alignment and grid offset), the plateau height $n\ge n_*=1$ and the shell $U=T\setminus\Hboxp$. The \emph{shell model} at shell mass $\varkappa\in[0,\infty)$ is King's walled torus model of \Cref{rot:def:tori}(c) with the wall $n\psi+\varkappa\mathbf 1_U$: its matter operator is $-\Delta_A+m^2+\dm_\eps+2n\psi+2\varkappa\mathbf 1_U$, and $\varkappa\mathbf 1_U$ is kept inside the matter covariance of every block-spin step. King's block-spin flow \cite[\S3]{K1} is run on $\Lk$, $\eps=\eps_k=5^{-k}$, through $r$ steps down to the \emph{stopping scale} $\eta_0\defeq5^{r-k}\in5^{-\N}$, at which the plateau is massive in block units (no quantitative heaviness of the plateau is used; only $n\ge1$ enters). The shell mass in block units at the stopping scale is $w_r\defeq\varkappa\eta_0^2$; the \emph{frozen regime} is $w_r\ge c_0$ for a fixed $c_0\in(0,c_0^*]$, and the threshold of \Cref{rot:hyp:HS} is
\[
\kz=\varkappa_0(\eta_0)\defeq\hat c_0/\eta_0^2,\qquad\hat c_0\defeq\max(c_0,1).
\]
The output of the flow is the effective action $S^{(r),\varkappa}(A_r,\phi_r;zf,0)$ on the block fields of the coarse lattice $5^r\eps\,\Z^3/P_{J'}$, in King's unit scale, with the curvature source $zf$ carried as in \cite[(3.20)]{K1}, and the small-field indicator $\chi_r$ of \cite[(3.2)]{K1}; the shell model carries in addition two tests at the stopping scale, a \emph{near-shell test} $\tau_r$ on the block fields of $U$ and a \emph{first-layer test} $\lambda_r$ on the block fields of the first light layer of $\Hboxp$, both with thresholds that depend on $w_r$ through increasing functions, and a weakened gradient family $\chi_r^\flat\ge\chi_r$ with no clause on the bonds inside $U$ or crossing $\partial U$. The stopping-scale integrals are
\[
\Ish_\varkappa(z)\defeq\int(dA_r)(d\phi_r)\,\chi_r\tau_r\lambda_r\,e^{-S^{(r),\varkappa}},\qquad
\Ifull_\varkappa(z)\defeq\int(dA_r)(d\phi_r)\,\chi_r^\flat\tau_r\lambda_r\,e^{-S^{(r),\varkappa}}\ \ge\ \Ish_\varkappa(z).
\]
For $\varkappa\ge\kz$ the thresholds of $\tau_r$, $\lambda_r$ are dominated by the $\varkappa$-free caps $C_7p_r$ on $U$ and $C_8p_r$ on the first layer ($p_r=p(\eta_0)$ King's small-field function), and the \emph{capped} integral $\Icap_\varkappa(z)$ of \Cref{rot:hyp:HS} is $\Ifull_\varkappa$ with $\tau_r$, $\lambda_r$ replaced by the indicators of these caps. Under the standing smallness $\eta_0\le5/(\lambda C_7^4)$ the cap on $U$ implies the block-sup clause of $\chi_r^\flat$ there, so the capped test set is a product: with $b\defeq(A_r,\phi_r|_{\Hboxp})$ the \emph{exterior block data} and $\phi_U\defeq\phi_r|_U$,
\[
\chicap\defeq\{b:\ \chi_r^\flat\lambda^{\rm cap}_r(b)=1\},\qquad
\Scap_\varkappa(b;z)\defeq-\ln\int\prod_{y\in U^{(r)}}d\phi_r(y)\ \mathbf 1\{|\phi_r|\le C_7p_r\text{ on }U^{(r)}\}\,e^{-S^{(r),\varkappa}(A_r,\phi_r;zf,0)},
\]
so that $\Icap_\varkappa(z)=\int_{\chicap}e^{-\Scap_\varkappa(b;z)}\,db$ exactly: this is the measure space $(\tilde\chi,db)$ and the reduced small-field action of \Cref{rot:hyp:HS}, and the test set $\chicap$ is one and the same for every $\varkappa\ge\kz$, which is the threshold clause of that theorem. Every statement below is uniform in $\eps\le\eps_0(\eta_0)$, in $\varkappa$ on the range stated, in $z\in[-1,1]$, and holds for the model with the free-boundary counterterm $\dm_\eps$ of \Cref{def:s02:acube}: the torus-convention counterterm of \cite{K1} differs from it by a number $\Delta_T(\eps)$ bounded uniformly in $\eps$ (\Cref{rot:prop:CT}(b) above), and every bound of the shell analysis holds with constants depending on $\sup_\eps|\Delta_T(\eps)|$ as well.

\begin{proposition}[the two-sided bound with the shell inside]\label{rot:prop:SW1}
Let $L=5$, the couplings $e,\lambda,m,\lmu$ in King's range, $(\Hboxp,T)$ block-aligned with the grid offset as in \Cref{rot:hyp:HS}, $n\ge1$, and $f$ as in \Cref{rot:hyp:HS}. There are $\eta_1>0$, $\sigma>0$, $c>0$ and, for every $\eta_0\in5^{-\N}$ with $\eta_0\le\eta_1$, a spacing bound $\eps_0(\eta_0)>0$ and a constant $C<\infty$ depending on $(\lambda,e,m,\lmu,|T|,n,f,\Hboxp,c_0)$ but not on $\eta_0$, $\varkappa$ or $\eps$, such that for all $\eps\le\eps_0(\eta_0)$ and $z\in[-1,1]$, with
\[
\sfa(\eta_0)\defeq C\eta_0^{\sigma}|T|,\qquad\sfb(\eta_0)\defeq e^{-c\,p(\eta_0)^2+C|T|},
\]
\begin{align*}
Z_\varkappa(zf,0)&\ \ge\ e^{-\sfa}\,\Ish_\varkappa(z)&&(\varkappa\in[0,\infty)),\\
Z_\varkappa(zf,0)&\ \ge\ e^{-\sfa}\,\Icap_\varkappa(z)&&(\varkappa\in[\kz,\infty)),\\
Z_\varkappa(zf,0)&\ \le\ e^{\sfa}\,\Ifull_\varkappa(z)+\sfb\,Z_\varkappa(0,0)\ \le\ e^{\sfa}\,\Icap_\varkappa(z)+\sfb\,Z_\varkappa(0,0)&&(\varkappa\in[0,\infty)).
\end{align*}
In particular $\sfa,\sfb\to0$ as $\eta_0\to0$ at fixed $(n,|T|)$, uniformly in $\varkappa$ and $\eps\le\eps_0(\eta_0)$, and clause \textup{(i)} of \Cref{rot:hyp:HS} holds for $\varkappa\ge\kz$ with the capped integral.
\end{proposition}

\begin{proof}
The two lines are King's Theorem~3.1 \cite[(3.3)--(3.4)]{K1}, that is Ba{\l}aban's lower bound \cite[\S\S3--4]{Ba1} and upper bound \cite[\S\S2--3]{Ba2} with the renormalisation of \cite{Ba3}, re-run for the shell model, at block factor $5$ (\Cref{thm:a06:BALA}) and with the remainder of the upper line measured against $Z_\varkappa(0,0)$. The proof has four layers, each uniform in $\varkappa\in[0,\infty]$.

\emph{(1) Exact algebra.} The block-spin transformation is a Gaussian identity valid for every positive matter covariance: with $\varkappa\mathbf 1_U$ a non-negative diagonal operator, every object of \cite[(2.13)--(2.17), (3.14)]{K1} is defined, the composition law of the steps, the scale decomposition of the fluctuation covariance into slices, the minimiser identity for the background $\phi^{(j)}_\varkappa=a_j(M^\varkappa_j)^{-1}Q_j(A)^*\phi_j$ and the Schur-complement monotonicity in $\varkappa$ hold verbatim, and every scalar line of the flow is a leg of block data (the slice covariances are $\ell\,C^{(j),\varkappa}\ell^\dagger$ with legs $\ell$). A two-line positivity argument inserts the near-shell test at every scale at no cost.

\emph{(2) Operator bounds.} King's Theorem~3.3 \cite[(3.5)--(3.8)]{K1} for the shell model in the forms consumed: exponential decay, sup and first-difference bounds for the step operators $G_j$, $\Delta^{(j)}$, $C^{(j)}$ everywhere, uniformly in $\varkappa$ (coercivity, since $\varkappa\mathbf 1_U\ge0$ adds to the quadratic form; $\ell^2$ decay by conjugation, which $\varkappa\mathbf 1_U$ commutes with; Loewner monotonicity for the unit-lattice kernels; Kato's inequality with the free Laplacian for the fine sup bounds); the derivative clause for the background of small-field data; H\"older bounds on the light side; and on the heavy side the amplitude bounds of a barrier argument (a constant supersolution of mass $\sqrt{w_j}$ in block units), with the first-difference bounds at the heavy points paired with the amplitude of the other leg of the pair. The pointwise difference bounds of the line kernels, $|D_AG(x,y)|\le C|x-y|_\eta^{-2}e^{-c|x-y|}$ and, with one difference at each end of one line, $|D^x_AD^y_AG(x,y)|\le C|x-y|_\eta^{-3}e^{-c|x-y|}$, uniformly in $\varkappa$ on both sides of $\partial U$, follow from the scale decomposition of (1), one difference per variable carried by its own leg, the leg bounds and a geometric sum over the scales. On the Neumann sub-domains of Ba{\l}aban's splitting the first-difference and H\"older clauses are used at points at distance at least Ba{\l}aban's localisation radius from the complement, where they hold; the convexity of $\partial U=\partial\Hboxp$ enters here, since at a re-entrant edge the first-difference bound fails uniformly in the spacing.

\emph{(3) Perturbative bounds.} Ba{\l}aban's renormalisation bound \cite[Prop.~1, (1.33)]{Ba3} for the renormalised families of \cite[\S3.5]{K1} of order at most $13$ with at least one external field, with the shell model's propagators and legs, in Ba{\l}aban's product form with a $\varkappa$-adapted norm on the scalar external fields, and the restricted form of King's bound \cite[(3.38)]{K1} on the field-dependent families; the Taylor step at a divergent two-scalar-leg subgraph acts on the photon legs, whose norms do not see the shell mass, with the remainder controlled by the current two-point kernel bound $|\Pi^\varkappa_j(x,y)|\le C|x-y|^{-4}$ from the difference bounds of (2); the two-photon-leg and mixed classes are combined by the Ward--Takahashi identities, which hold for the shell model since $\varkappa\mathbf 1_U$ is gauge invariant, used for full step-$j$ families and never for a fixed ordering of the lines; a mismatch vertex $\mathsf M$ (the difference of the shell-adapted and the bulk subtraction inside a renormalised pair) is a two-leg vertex of bounded coefficient. The block-data form of the bound, with King's legs $\phi^{(r)}_\varkappa$ and the small-field data $|\phi_r|\le p$, $|\partial A_r|\le p$, $|D_{A_r}\phi_r|\le p$ of \cite[(3.2)]{K1}, is derived from the product form on the photon legs by a local gauge transformation that removes the constant part of the background before the field itself is bounded, King's data satisfying $|A_r|\le p\,(\lmu\ell_{r-1})^{-1}$ but no uniform bound; the same reduction gives the light-side first moment of the photon two-leg kernel, summed against the background and its first differences, which the near-face analysis of \cite[\S5.2]{K2} transplanted to the shell model consumes. \moving{proof of the bound, uniform in the shell mass, on the first-order term in the background field of the determinant of the step fluctuation Gaussian on the Neumann sub-domains of Ba{\l}aban's splitting, where the large-field analysis of the upper bound uses it: being completed}

\emph{(4) The two bounds.} The lower bound is \cite[\S\S3--4]{Ba1}: restriction to the small-field region by Jensen's inequality, the inductive hypothesis on King's Appendix class \cite[(A.1)--(A.8)]{K1} for the shell model (with a background-gradient and a block-mismatch clause, which compose from step to step), the cumulant bound for the step Gaussian with $\varkappa$-free constants (\Cref{rot:thm:R1}), and the final step \cite[(3.66)--(3.69)]{Ba1}; the restriction to the capped set at the stopping scale is a positivity step inside the class, which gives the second line for $\varkappa\ge\kz$. The upper bound is \cite[\S\S2--3]{Ba2}: the large/small-field split at every scale, with the near-shell test paid at every scale and the first-layer test at the stopping scale, the sum over large-field configurations $\prod_j\sum_\Lambda e^{-p(\ell_j)|\Lambda|}\le e^{C|T|}$, the propagator-domain comparisons of \cite[Prop.~2.6]{Ba2} in the form \Cref{hyp:k01:BALint} with the regrouping \Cref{hyp:k01:RENorb}, and the continuation of the large-field terms past the stopping scale, which under block alignment runs the shell model at the scales below $\eta_0$; the remainder is $\sfb\,Z_\varkappa(0,0)$ in normalisation-free form by the ultraviolet-stability gap, as in \eqref{eq:k01:H0rel}. The upper capped line follows from $\Ifull_\varkappa\le\Icap_\varkappa$ (the cap tests dominate $\tau_r$, $\lambda_r$ pointwise as indicators and the integrand is positive); the lower capped line is the positivity step inside the class just described, the capped test set lying in King's Appendix class for $\varkappa\ge\kz$.
\end{proof}

\begin{proposition}[the $\varkappa$-Cauchy property of the reduced action]\label{rot:prop:SW3}
In the setting of \Cref{rot:prop:SW1} fix $\alpha_0\in(0,\frac12]$. There is $\eta_2\le\eta_1$ such that for every $\eta_0\in5^{-\N}$ with $\eta_0\le\eta_2$ there is $C_{\mathfrak c}(\eta_0)<\infty$, depending on $(T,\Hboxp,n,\lambda,e,m,\lmu,\alpha_0,c_0,\eta_0)$ and on $f$ but not on $\eps$, $\varkappa$, $z$, such that for all $\eps\le\eps_0(\eta_0)$, all $\kz\le\varkappa\le\varkappa'<\infty$ and all $z\in[-1,1]$,
\[
\sup_{b,b'\in\chicap}\big|\Delta_z(b)-\Delta_0(b')\big|\ \le\ 2R(\varkappa,\varkappa')\ \le\ 2\,r(\varkappa;\eta_0),\qquad
R(\varkappa,\varkappa')\defeq\int_{\varkappa\eta_0^2}^{\varkappa'\eta_0^2}\theta(w)\,dw+\tau(\varkappa),
\]
\[
\theta(w)\defeq C_{\mathfrak c}\,w^{-1}w^{-(1-\alpha_0)/2}(1+\log w)^2,\qquad
\tau(\varkappa)\defeq C_{\mathfrak c}\,(\varkappa\eta_0^2)^{-\frac12+\frac{\alpha_0}4},\qquad
r(K;\eta_0)\defeq C'_{\mathfrak c}\,(K\eta_0^2)^{-(1-\alpha_0)/2}\big(1+\log(K\eta_0^2)\big)^2,
\]
with $\Delta_z(b)=\Scap_{\varkappa'}(b;z)-\Scap_\varkappa(b;z)$ as in \Cref{rot:hyp:HS}; $r(K;\eta_0)\to0$ as $K\to\infty$ at fixed $\eta_0$. This is clause \textup{(ii)} of \Cref{rot:hyp:HS}.
\end{proposition}

\begin{proof}
The statement is derived from a bound on King's full explicit action in all block fields: there is a field-free function $\varkappa\mapsto c_\varkappa$ on $[\kz,\infty)$, independent of $z$, with
\[
\sup_{(A_r,\phi_r)\in\supp(\chi^\flat_r\tau^{\rm cap}_r\lambda^{\rm cap}_r)}\big|S^{(r),\varkappa}-S^{(r),\varkappa'}-(c_\varkappa-c_{\varkappa'})\big|\ \le\ R(\varkappa,\varkappa'),
\]
and on the capped set, whose $U$-section is the same cube $\{|\phi_U|\le C_7p_r\}$ for every exterior datum $b$ and every $\varkappa$, the reduced actions inherit the bound with the same $c_\varkappa$ (the logarithm of an integral of $e^{-S}$ over a fixed set moves by at most the sup of the change of $S$); taking the mixed difference removes $c_\varkappa$. The full-field bound is proved in derivative form: $\partial_\varkappa S^{(r),\varkappa}$ is King's sum of renormalised graphs with one insertion of $\partial_\varkappa$ of the matter operator, that is of $2\mathbf 1_U$ on a line, plus the $\varkappa$-derivative of the block-free part of the background determinant. For an insertion on a line with a scalar connection to the rest of the graph the heavy matter on $U$ pays $w_j^{-1/2}$ per deep end at every frozen scale (the barrier bounds of layer (2) of the proof of \Cref{rot:prop:SW1}, with the one-loop $z^2$-term and its positions in $U$ treated in both regimes); for an insertion whose matter component through it is closed and connects only through photon lines, which do not see the shell mass, the members with the same full lines are grouped into one class, one gauge identity is applied to the class, and only then is it expanded into King's members with the background rebased for each, the base layer being a first-moment bound, uniform in $\varkappa$, for the two-photon kernel built from full lines. In the units of the stopping scale, with $\hat\varkappa=\varkappa\eta_0^2$, the sum over the members less its field-free part is at most $\theta(\hat\varkappa)$ per unit of $\hat\varkappa$, integrable at infinity; the H\"older remainders of the Taylor step for lines used as legs and the lattice regime of the scale sum give the exponent $(1-\alpha_0)/2$. The block-free part of the determinant has the monotone tail $0\le\Phi^{\rm bf}(\varkappa)-\Phi^{\rm bf}(\infty)\le\tau(\varkappa)$, with the bracket for the loops winding through $U$ (whose contribution is positive even at zero background curvature) bounded, under block alignment, by $C\hat\varkappa^{-1/2}$. Integrating in $\varkappa$ gives $R(\varkappa,\varkappa')$, and the sum over the dyadic intervals of $\hat\varkappa$ gives $r(K;\eta_0)$. The perturbative inputs are those of layer (3) of the proof of \Cref{rot:prop:SW1}, the restricted form of \cite[(3.38)]{K1} being re-run for the shell vertex set; the block-data form is consumed at the step marked there.
\end{proof}

\begin{proof}[Proof of \Cref{rot:hyp:HS}]
Take $\eta_1$ of the theorem to be the $\eta_2$ of \Cref{rot:prop:SW3}, $\sfa,\sfb$ and $\eps_0(\eta_0)$ as in \Cref{rot:prop:SW1}, $\varkappa_0=\kz$, $r(\cdot;\eta_0)$ as in \Cref{rot:prop:SW3}, and $(\tilde\chi,db)=(\chicap,db)$ with $\tilde S_\varkappa=\Scap_\varkappa$. Clause (i) is the second and third lines of \Cref{rot:prop:SW1} for $\varkappa\ge\kz$; clause (ii) is \Cref{rot:prop:SW3}; the test set does not depend on $\varkappa$ on $[\kz,\infty)$, as stated.
\end{proof}


In words. King's block-spin flow \cite[\S3]{K1} is run for the model in which $\varkappa\mathbf 1_U$ is part of the matter covariance at every scale, down to a stopping scale $\eta_0$ at which the plateau is massive in block units; $\varkappa_0(\eta_0)$, of the order of $\eta_0^{-2}$, is the mass above which the shell is heavy in block units at that scale. $\Icap_\varkappa(z)$ is the stopping-scale integral of the flow over a capped set $\tilde\chi$ of block fields, the \emph{test set}, and $\tilde S_\varkappa(b;z)$ is the reduced small-field action at the block field $b$, shell mass $\varkappa$ and source coefficient $z$, both for the model at the spacing $\eps_k$: the flow from $\eps_k$ to $\eta_0$ depends on $\eps_k$, and so do the test set and the reduced action, whereas the constants $\sfa,\sfb,\varkappa_0,\eps_0$ and the function $r$ do not. Clause (i) is the two-sided bound of King's and Ba{\l}aban's type (small-field integral below, small-field integral plus a large-field remainder above) for the model with the shell mass inside the covariance, uniformly in $\varkappa\ge\varkappa_0$ and in the spacing; clause (ii) says that the dependence of the reduced action on the shell mass is, up to a quantity that tends to $0$ as the shell gets heavier, the same at every block field and every source coefficient. That the test set is one and the same for all $\varkappa\ge\varkappa_0$ is where the threshold enters: no term from a change of representation of the stopping-scale integral arises between $\varkappa$ and $\varkappa'$, and (ii) compares two integrands on one set. The proof is \Cref{rot:prop:SW1,rot:prop:SW3}; \Cref{rot:rem:behindHS} lists its inputs beyond the printed literature.

\begin{theorem}[the heavy-shell estimate from the shell bounds]\label{rot:thm:shell}
Let $g=0$ and let $(\fs,\chi)$ be a two-level configuration with the exterior shell $U=T\setminus\Hboxp$, plateau height $n\ge n_*$ and $\Hboxp$, $T$ satisfying the block-alignment and grid-offset clauses of \Cref{rot:hyp:HS}, and let $f$ be a real $2$-form in $C^\infty_c$ supported in $\Qb b$. Then for every $\eta>0$ there are a stopping scale $\eta_0$ and a shell mass $K<\infty$ with
\[
\limsup_{k\to\infty}\tail_k(K)\le\eta ;
\]
in particular \REz\ holds for $(f,0)$ and $\eta$ at every such configuration with $\tauGB(\fs)\le\eta$, with this $K$. Consequently the closing estimate \TIz\ holds (\Cref{rot:thm:REStoTIGHT}, with \Cref{rot:hyp:K,rot:hyp:DS}), with the single-level configuration $(J',\psi,n+K)$ in the class of \Cref{rot:rem:class}.
\end{theorem}

\begin{proof}
Fix $(\fs,\chi)$, the source $f$ and $g=0$. For $\varkappa\in[0,\infty]$ and $z\in\C$ put $\Psi_\varkappa(z)\defeq Z_\varkappa(zf,0)/Z_\varkappa(0,0)=\langle e^{zF_{\Lk}(f)}\rangle_\varkappa$, where for $\varkappa=\infty$ the state is the walled torus state with the matter killed on $\Dchi=U$. The shell mass $\varkappa\mathbf 1_U$ is the second level $\varkappa\chi$ of the wall, so by \Cref{rot:def:twolevel}, at $g=0$, $\Psi_K(\mathrm i)=\Psik(n,K)$, $\Psi_\infty(\mathrm i)=\Psik(n,\infty)$ and $\tail_k(K)=\Psi_K(\mathrm i)-\Psi_\infty(\mathrm i)$. For $\varkappa\ge\varkappa_0$ and $z\in[-1,1]$ put $\cJ_\varkappa(z)\defeq\Icap_\varkappa(z)/\Icap_\varkappa(0)$, with the test set and the reduced action of \HS\ at the spacing $\eps_k$ under consideration. Every constant below is uniform in $k$ with $\eps_k\le\eps_0(\eta_0)$.

\emph{Step 0: the disc bound.} At fixed spacing $z\mapsto\Psi_\varkappa(z)$ is entire, a finite-dimensional integral of $e^{zF(f)}$ against a measure with Gaussian photon tails. For $|z|\le2$, $|e^{zF(f)}|\le e^{2F(f)}+e^{-2F(f)}$. Each state $\langle\cdot\rangle_\varkappa$, $\varkappa\in[0,\infty]$, is a BFS measure whose Gaussian is the periodic massive photon Gaussian of the torus, with covariance $C$; the wall, the shell mass and the killing act in the single-site measures only. So \Cref{lem:s02:GD} with $s=\pm2$ and $v$ the vector with $\ip vA=F_{\Lk}(f)$ gives $\langle e^{\pm2F(f)}\rangle_\varkappa\le e^{2\ip v{Cv}}$, hence
\[
|\Psi_\varkappa(z)|\le\Mdisc\defeq2e^{2\ip v{Cv}}\qquad(|z|\le2,\ \varkappa\in[0,\infty]).
\]
Here $\ip v{Cv}$, the variance of $F(f)$ under the free massive photon on the torus, is bounded uniformly in the spacing for smooth $f$: the periodic photon form is at least $\lmu^2$ times the $\ell^2$ norm, so $\ip v{Cv}\le\lmu^{-2}\norm{\dstar_kf}^2_{2,k}$, a Riemann sum for $\lmu^{-2}\norm{\dstar f}^2_{L^2}$ (as in \Cref{rot:cor:walled}). At $g=0$ no matter moment enters the disc bound. For real $z$, $\Psi_\varkappa(z)\to\Psi_\infty(z)$ as $\varkappa\to\infty$ at fixed spacing, by the substitution and dominated convergence of the proof of \Cref{lem:s02:Kprime} applied with the set $U$, the rate $\chi\equiv1$ on $U$ and $\Theta=e^{zF(f)}$, which is of the form admitted there ($|\Theta|=e^{\ell(A)}$ with $\ell=zF(f)$ real linear), the factor $e^{zF(f)}$ being integrable against the photon Gaussian by the same bound.

\emph{Step 1: the sandwich.} Let $\eta_0\le\eta_1$ be so small that $\sfb(\eta_0)\le\frac12$, fix $k$ with $\eps_k\le\eps_0(\eta_0)$, and let $\varkappa\ge\varkappa_0$, $z\in[-1,1]$; write $\sfa=\sfa(\eta_0)$, $\sfb=\sfb(\eta_0)$. By \HS(i) at $z=0$, $e^{-\sfa}(1-\sfb)Z_\varkappa(0,0)\le\Icap_\varkappa(0)\le e^{\sfa}Z_\varkappa(0,0)$; by \HS(i) at $z$ and these two bounds,
\[
e^{-2\sfa}(1-\sfb)\,\cJ_\varkappa(z)\ \le\ \Psi_\varkappa(z)\ \le\ e^{2\sfa}\cJ_\varkappa(z)+\sfb .
\]
Hence $\cJ_\varkappa(z)\le2e^{2\sfa}\Mdisc=:\bar{\cJ}$ and $|\Psi_\varkappa(z)-\cJ_\varkappa(z)|\le\bar{\cJ}(e^{2\sfa}-e^{-2\sfa}+\sfb)+\sfb=: m_1(\eta_0)$, and $m_1(\eta_0)\to0$ as $\eta_0\to0$ because $\sfa,\sfb\to0$.

\emph{Step 2: the $\varkappa$-Cauchy property.} Let $\varkappa_0\le\varkappa\le\varkappa'<\infty$ and $z\in[-1,1]$. Since at this $k$ the test set $\tilde\chi$ is the same for $\varkappa$ and $\varkappa'$,
\[
\frac{\cJ_\varkappa(z)}{\cJ_{\varkappa'}(z)}=\frac{\iint_{\tilde\chi\times\tilde\chi}e^{-\tilde S_\varkappa(b;z)-\tilde S_{\varkappa'}(b';0)}\,db\,db'}
{\iint_{\tilde\chi\times\tilde\chi}e^{-\tilde S_{\varkappa'}(b;z)-\tilde S_\varkappa(b';0)}\,db\,db'},
\]
and the two integrands differ pointwise by the factor $e^{\Delta_z(b)-\Delta_0(b')}$, which lies in $[e^{-2r},e^{2r}]$ with $r\defeq r(\varkappa;\eta_0)$ by \HS(ii). So $\cJ_\varkappa(z)/\cJ_{\varkappa'}(z)\in[e^{-2r},e^{2r}]$ and $|\cJ_\varkappa(z)-\cJ_{\varkappa'}(z)|\le\bar{\cJ}(e^{2r}-1)$. Put $\bar r(K;\eta_0)\defeq\sup_{\varkappa\ge K}r(\varkappa;\eta_0)$, which tends to $0$ as $K\to\infty$ at fixed $\eta_0$. With Step 1, for all $\varkappa'\ge\varkappa\ge K\ge\varkappa_0$,
\[
|\Psi_\varkappa(z)-\Psi_{\varkappa'}(z)|\ \le\ m(K,\eta_0)\defeq2m_1(\eta_0)+\bar{\cJ}\big(e^{2\bar r(K;\eta_0)}-1\big),\qquad z\in[-1,1],
\]
uniformly in $k$ with $\eps_k\le\eps_0(\eta_0)$.

\emph{Step 3.} Letting $\varkappa'\to\infty$ at fixed spacing (Step 0): $|\Psi_K(z)-\Psi_\infty(z)|\le m(K,\eta_0)$ for $z\in[-1,1]$ and $K\ge\varkappa_0$.

\emph{Step 4: two constants.} $\Xdiff\defeq\Psi_K-\Psi_\infty$ is entire, $|\Xdiff|\le2\Mdisc$ on $\{|z|\le2\}$ and $|\Xdiff|\le m(K,\eta_0)$ on $[-1,1]$. If $m(K,\eta_0)=0$, then $\Xdiff$ vanishes on $[-1,1]$ and so identically. Otherwise $\log|\Xdiff|$ is subharmonic on the domain $\Omega\defeq\{|z|<2\}\setminus[-1,1]$, bounded above by $\log(2\Mdisc)$ there, with upper limit at most $\log m(K,\eta_0)$ at every point of the slit; by the two-constants theorem \cite[\S4.3]{Ran95} it is dominated on $\Omega$ by the harmonic function with boundary values $\log(2\Mdisc)$ on $|z|=2$ and $\log m(K,\eta_0)$ on the slit. At the point $\mathrm i\in\Omega$,
\[
|\Xdiff(\mathrm i)|\le(2\Mdisc)^{1-\omega}\,m(K,\eta_0)^{\omega},
\]
where $\omega\in(0,1)$ is the harmonic measure of the slit $[-1,1]$ at $\mathrm i$ in $\Omega$, an absolute constant.

\emph{Step 5: the real point.} The charge conjugation $(A,\vartheta)\mapsto(-A,-\vartheta)$ preserves the action, the wall, the shell mass and the measure (\Cref{rot:def:tori}), and reverses the sign of $F(f)$; so $\Psi_\varkappa(-z)=\Psi_\varkappa(z)$, and since $\Psi_\varkappa(\bar z)=\overline{\Psi_\varkappa(z)}$, the value $\Psi_\varkappa(\mathrm i)$ is real. Thus $\Xdiff(\mathrm i)=\tail_k(K)\ge0$ (\Cref{rot:lem:tailprop}(a)) and
\[
0\le\tail_k(K)\le(2\Mdisc)^{1-\omega}\,m(K,\eta_0)^{\omega}\qquad\text{for every }\eps_k\le\eps_0(\eta_0).
\]

\emph{Choice of the parameters.} Given $\eta>0$: take $\delta>0$ with $(2\Mdisc)^{1-\omega}\delta^{\omega}\le\eta$; then $\eta_0\le\eta_1$ with $\sfb(\eta_0)\le\frac12$ and $2m_1(\eta_0)\le\delta/2$, which is possible because $\sfa,\sfb\to0$; then $K\ge\varkappa_0(\eta_0)$ with $\bar{\cJ}(e^{2\bar r(K;\eta_0)}-1)\le\delta/2$, which is possible because $r(K;\eta_0)\to0$ as $K\to\infty$ at fixed $\eta_0$; and only then $\eps_k\le\eps_0(\eta_0)$. Then $m(K,\eta_0)\le\delta$ and $\tail_k(K)\le\eta$ for all $\eps_k\le\eps_0(\eta_0)$, so $\limsup_k\tail_k(K)\le\eta$. This is the bracket of \REz\ at $(\fs,\chi)$ with this $K$, in the quantifier order of \Cref{rot:hyp:RES}: one configuration, one $K$, no rate, the limit along the spacings last; since $\tauGB(\fs)$ depends only on $\Hbox$ and the width $D(\fs)$, a configuration with $\tauGB(\fs)\le\eta$, with $b'$ in the residue class of the grid-offset condition and $T$ a King torus, so that $\Hboxp$ and $T$ satisfy the alignment clauses of \Cref{rot:hyp:HS} (every interval of length $M'$ final blocks contains an admissible $b'$), is available for every $b$, $f$ and $\eta$, and \REz\ holds. Since $\tail_k$ is nonincreasing in $K$ (\Cref{rot:lem:tailprop}(a)), $K$ may be enlarged at will (\Cref{rot:rem:Eproof}(iii)), which places $(J',\psi,n+K)$ in the class of \Cref{rot:rem:class}. The proof of \Cref{rot:thm:REStoTIGHT} uses the dead-shell bound \eqref{rot:eq:deadbound} at this configuration, which is \DSh\ (\Cref{rot:hyp:DS}), and the existence of the walled value of $(J',\psi,n+K)$, which is part of \Kz\ (\Cref{rot:hyp:K}), together with King's printed state (\EC); it then gives \TIz.
\end{proof}

\begin{proof}[Proof of \Cref{rot:thm:main}]
\Cref{rot:thm:shell} gives, from the heavy-shell estimates \HS\ (\Cref{rot:hyp:HS}), the decoupling \REz\ at one two-level configuration for each $b$, $f$ and $\eta$, and its final clause, \Cref{rot:thm:REStoTIGHT} with the torus comparison \Kz\ (\Cref{rot:hyp:K}) and the dead-shell chain \DSh\ (\Cref{rot:hyp:DS}), gives the closing estimate \TIz. With \EC\ (\Cref{hyp:s03:EC}), \Kz\ and \TIz, \Cref{rot:thm:curv} gives the invariance of $\ZKinf(\cdot,0)$ under $O(3)$ for all complex coefficients, and \Cref{cor:s06:TK} adds the translations.
\end{proof}

\begin{remark}[the inputs of the proof of \HS]\label{rot:rem:behindHS}
The two clauses of \HS\ are proved by a renormalisation-group analysis of King's walled torus model \cite{K1,K2} with the shell mass kept inside the matter covariance at every scale, run down to a stopping scale at which the plateau is massive in block units. Its ingredients beyond the printed literature, by content: Ba{\l}aban's asserted block-factor-$5$ estimates re-run for the shell model; a frozen-scale sub-domain bound in scale-free form; the sub-domain form of the perturbative bound with a tangentially localised amplitude; pointwise first- and mixed second-difference bounds on the line kernels, on the torus and on the Neumann sub-domains used, uniformly in the shell mass and on both sides of $\partial U$; the light-side regularity estimates on sub-domains and a localisation lemma for the local inverses near the edges and corners of the source box; the tail of the block-free determinant factor with its bracket for the loops winding through $U$; heavy-kernel bounds for the matter operator with the shell mass; the $\varkappa$-derivative bound for clusters closed in matter, through a gauge identity for full lines; the restricted renormalisation theorem for the shell vertex set; a first-moment bound, uniform in $\varkappa$, for the two-photon kernel built from full lines; King's near-face analysis \cite[\S5.2]{K2} transplanted to the shell model; Ba{\l}aban's all-orders enumeration \cite[\S\S2--3]{Ba3} with one inserted vertex; King's Appendix class \cite[(A.1)--(A.8)]{K1} for the shell model. The counterterm offset between the torus and the free-boundary convention enters as a bounded constant. Of the printed literature it uses, besides the theorems as printed, the statements named in \Cref{rot:rem:behindK} together with \cite[Lemma~2.7, (2.113)]{Ba2}; the shell analysis shares with the torus comparison and with the existence theorem (\cref{sec:k01}) Ba{\l}aban's same-scale estimates at block factor $5$, their reading at a scalar wall with unwalled photon lines, and the orbit regrouping of the cut graph families. Several forms in which such shell estimates might be posed are false and are not used: a sub-domain clause with an absolute constant and a rate in the penetration depth, which fails when the heaviness threshold of the shell in block units is small (the scale-free form is used); the determinant tail without the winding bracket and with a right side proportional to the square of the background curvature, which fails on the torus (a constant background has zero curvature while the winding loops contribute a positive tail); a Yukawa form of the heavy-pair clause, which fails for pairs in the first layer; and a supremum bound for the full covariance next to $\partial U$ in the lattice regime. \Cref{rot:thm:shell} uses the determinant tail only through a bound on it that tends to $0$ as $\varkappa\to\infty$, so a change of its exponent or a logarithmic factor leaves the theorem unchanged. A photon Ward identity for each fixed ordering of King's graphs would not supply the $\varkappa$-derivative bound, since that identity is false; the identities that hold are those for full lines and for Ba{\l}aban's resummed organisation, which keeps a convergent remainder, and in its place the bound rests on the first-moment bound for the two-photon kernel. \DSh\ (\Cref{rot:hyp:DS}) carries the three steps of \Cref{rot:thm:SHtoTIGHT}, whose proof is written for a shell between two aligned cubes, to the geometry with one convex face.
\end{remark}

\subsection{The continuum form}

This subsection is not on the chain of \Cref{rot:thm:main}: it records an equivalent form of \REa, which is not used, and the existence statement \Einf\ introduced here is used by nothing. Both sides of the bracket in \REa\ have continuum limits, the first by \Ka\ and the second under one further existence statement; in those terms the property says that two continuum objects coincide up to $\eta$. Throughout, $(\fs,\chi)$ is a two-level configuration for $\Qb b$ in the sense of \Cref{rot:def:shell} (smooth profile $\chi$, shell $U=C_L\setminus\Hbox$), and $(f,g)\in\cK$ is supported in $\Qb b$.

\begin{proposition}[the dead-shell state between two cube states]\label{rot:prop:deadcube}
Let $B_k\defeq e^{-a_k\sS_{\Lk}(g)}\langle\Theta(\mathrm if,g)\rangle$ in the Neumann--Dirichlet state on $C_L\cap\Lk$ (Neumann photon on $C_L$, Dirichlet matter outside $C_L$) whose single-site measures are those of the dead-shell walled torus: the wall $n\psi$, and the matter killed on $\Dchi\cap C_L$. Then at every $k$ with $\eps_k$ small
\[
B_k\;\le\;\Psik(n,\infty)\;\le\;B_k+Z^a_{k,L}(0,2g)^{1/2}\big(e^{\tauGB(\fs)}-1\big)^{1/2},\qquad B_k\le Z^a_{k,L}(\mathrm if,g).
\]
\end{proposition}

\begin{proof}
The lower bound is \Cref{rot:lem:cutwall} applied to the dead-shell torus state and $X=C_L\cap\Lk$ ($X$ and its neighbours lie in one fundamental cell, as $\Nr1{\bar C_L}\subset(-25MJ',25MJ')^3$): the only couplings between $X$ and its complement are photon gradient terms, since every site of $C_L$ adjacent to the complement is dead, and deleting them lowers the expectation and factorises the measure, whose $X$-factor is the state defining $B_k$. The upper bound is Steps~1--2 of the proof of \Cref{rot:thm:SHtoTIGHT}, since the $N$-measure there is exactly the state defining $B_k$, and $e^{-2a_k\sS_{\Lk}(g)}\langle\Theta(0,2g)\rangle_N\le Z^a_{k,L}(0,2g)$ by its Step~3. The last inequality is that Step~3.
\end{proof}

\begin{theorem}[the continuum form]\label{rot:thm:cont}
Assume \Ka\ if $g\ne0$ (at $g=0$ the torus comparison is \Cref{rot:hyp:K}); for the sandwich in \textup{(d)} at $g\ne0$ assume also $\Wadm a$.
\begin{enumerate}[label=\textup{(\alph*)},leftmargin=*]
\item For every $K\in[0,\infty)$ the limit $\Psio(n,K)\defeq\lim_k\Psik(n,K)$ exists, equals the walled value $\Wv^a_{(J',n\psi+K\chi)}(f,g)$, and is nonincreasing in $K$; so $\Psio(n,\infty^-)\defeq\lim_{K\to\infty}\Psio(n,K)$ exists.
\item $\liminf_k\tail_k(K)=\Psio(n,K)-\limsup_k\Psik(n,\infty)$ for every $K$, and $\limsup_k\Psik(n,\infty)\le\Psio(n,\infty^-)$.
\item \REa\ holds at $(\fs,\chi)$ with tolerance $\eta$ for some $K$ if and only if
\[
\Psio(n,\infty^-)=\lim_{K\to\infty}\lim_{k\to\infty}\Psik(n,K)\;\le\;\limsup_{k\to\infty}\Psik(n,\infty)+\eta .
\]
\item If the \emph{dead-shell torus value} $\Psio(n,\infty)\defeq\lim_k\Psik(n,\infty)$ exists, then \textup{(c)} reads $\lim_{K\to\infty}\Psio(n,K)\le\Psio(n,\infty)+\eta$, the reverse inequality $\lim_K\Psio(n,K)\ge\Psio(n,\infty)$ always holds, and
\[
B^0\le\Psio(n,\infty)\le B^0+Z^a_\infty(0,2g)^{1/2}\big(e^{\tauGB(\fs)}-1\big)^{1/2}
\]
for every limit point $B^0$ of $(B_k)$; in particular $\Psio(n,\infty)\le Z^a_\infty(\mathrm if,g)+Z^a_\infty(0,2g)^{1/2}(e^{\tauGB(\fs)}-1)^{1/2}$. At $g=0$ the sandwich reads $B^0\le\Psio(n,\infty)\le B^0+(e^{\tauGB(\fs)}-1)^{1/2}$ and $\Psio(n,\infty)\le\ZKinf(\mathrm if,0)+(e^{\tauGB(\fs)}-1)^{1/2}$.
\end{enumerate}
\end{theorem}

\begin{proof}
(a) $(J',n\psi+K\chi)$ is a walled configuration in $\Aw{\Qb b}$ (its wall is smooth and vanishes on $\Nr1{\Qb b}$), so \Cref{rot:lem:K}(a) gives the limit, under \Kz\ at $g=0$; monotonicity in $K$ passes to the limit (\Cref{rot:lem:tailprop}(a)). (b) $\liminf_k(x_k-y_k)=\lim_kx_k-\limsup_ky_k$ when $x_k$ converges; and $\Psik(n,\infty)\le\Psik(n,K)$ for every $k$ and $K$. (c) By (b), \REa\ for some $K$ says $\Psio(n,K)\le\limsup_k\Psik(n,\infty)+\eta$ for some $K$, which by monotonicity in $K$ is the same as $\inf_K\Psio(n,K)\le\limsup_k\Psik(n,\infty)+\eta$. (d) The first two statements are (b), (c) with the limit in place of the $\limsup$; the sandwich is \Cref{rot:prop:deadcube} in the limit along a subsequence on which $B_k\to B^0$. At $g\ne0$, $Z^a_{k,L}\to Z^a_L\le Z^a_\infty$ on $\cK$ (\Cref{rot:thm:full}(i),(iii), under $\Wadm a$), with $(0,2g)\in\cK$. At $g=0$, $Z^a_{k,L}(0,0)=1$, and $B^0\le\limsup_kZ_{k,L}(\mathrm if)\le Z^K_{L+1}(\mathrm if)\le\ZKinf(\mathrm if,0)$ by \Cref{cor:s04:G2}(ii)($\gamma$) and the monotonicity of $J\mapsto\ZKJ(\mathrm if)$ (\Cref{thm:s05:CK}(b)); the support condition of that corollary holds since $\supp f\subset\Qb b\subset\Hbox$ and $\dist_\infty(\Hbox,\R^3\setminus C_L)\ge4$.
\end{proof}

\begin{definition}[\REcont: the continuum form]\label{rot:hyp:REScont}
\REcont\ is the property: for every $b>0$, $(f,g)\in\cK$ supported in $\Qb b$ and $\eta>0$ there is a two-level configuration $(\fs,\chi)$ for $\Qb b$ with $\tauGB(\fs)\le\eta$ for which the dead-shell torus value $\Psio(n,\infty)$ exists and
\[
\lim_{K\to\infty}\Psio(n,K)\;\le\;\Psio(n,\infty)+\eta .
\]
\end{definition}

By \Cref{rot:thm:cont}, \REcont\ implies \REa, and \REa\ implies \REcont\ whenever the dead-shell torus values of the configurations in question exist; the two are the same statement about the interchange of the limits $K\to\infty$ and $k\to\infty$, and in \REcont\ no lattice spacing appears. The existence of $\Psio(n,K)$ is covered by \Ka\ because the two-level wall $n\psi+K\chi$ is a smooth source; with the sharp profile $\mathbf 1_U$ it would not be. The existence of $\Psio(n,\infty)$, the continuum limit of a torus state with an interior region of Dirichlet matter, is the statement \Einf. It would be a theorem of King's type; King does not print it: the Dirichlet region is bounded by lattice hyperplanes only when $\partial\Hbox$ is, and it would have to come from the methods of \cite{K2} for Dirichlet matter on a cube, applied to a shell between two aligned cubes; the statement is not in \cite{K1,K2}. Neither \Einf\ nor \REcont\ is consumed by any result of this paper, and the lattice form \REa\ needs neither of them nor any limit beyond \Ka.

\subsection{The fixed-spacing statement and the shape of a counterexample}

\begin{proposition}[fixed spacing]\label{rot:prop:fixedeps}
For every two-level configuration, every $(f,g)\in\cK$ supported in $\Qb b$ and every $k$, $\tail_k(K)\to0$ as $K\to\infty$. Consequently, if \REa\ fails for some $(f,g)$, $\eta$, then for every two-level configuration $(\fs,\chi)$ for $\Qb b$ with $\tauGB(\fs)\le\eta$ and every $K$, $\tail_k(K)>\eta$ for all large $k$, while $\tail_k(K)\le\eta$ for $K\ge K_*(k)<\infty$; hence $K_*(k)\to\infty$ as $k\to\infty$.
\end{proposition}

\begin{proof}
The first statement is \Cref{rot:lem:tailprop}(a). The second is its contrapositive form: the failure of \REa\ at $(\fs,\chi)$ means $\liminf_k\tail_k(K)>\eta$ for every $K$, so $\tail_k(K)>\eta$ for all large $k$; and $K_*(k)\defeq\inf\{K:\tail_k(K)\le\eta\}$ is finite by the first statement. If $K_*(k)\le K_0$ along a subsequence, then $\tail_k(K_0)\le\eta$ along it, contradicting the previous sentence.
\end{proof}

So any failure of \REa\ would be carried by shell matter whose mass squared tends to infinity with $k$: a lattice-scale effect that disappears at every fixed cutoff.

\subsection{Why the wall, and why no soft argument}

\begin{proposition}[no common lattice hyperplanes]\label{rot:prop:forced}
Let $R=R_{\mu\nu}(\pm\thz)\in\Rgen$ and let $\rho$ be the third index.
\begin{enumerate}[label=\textup{(\alph*)},leftmargin=*]
\item There is no lattice $P\subset\R^3$ of rank $3$ with $RP=P$.
\item An affine plane is a lattice hyperplane of both $\Lk$ and $R^{-1}\Lk$ (that is, of the form $\{x_\alpha=c\}$ and of the form $\{(R^{-1}e_\beta)\cdot x=c'\}$) only if it is of the form $\{x_\rho=c\}$.
\item Consequently no bounded region of $\R^3$ has all its faces on lattice hyperplanes of both lattices; the only unbounded regions with that property are the slabs $\{c_1\le x_\rho\le c_2\}$. The bounded geometries to which both lattices descend are the tori $\R^3/P$ with $P\subset\Lk\cap R^{-1}\Lk$, for instance $P=50MJ\,\Z^3$; by \textup{(a)}, no such $P$ is $R$-invariant, so the rotated torus $R^{-1}\DJ$ is never the torus $\DJ$.
\end{enumerate}
\end{proposition}

\begin{proof}
(a) If $RP=P$, then in a $\Z$-basis of $P$ the map $R$ is an integer matrix, whose characteristic polynomial has integer coefficients; but the characteristic polynomial of $R$ is $(x-1)(x^2-\frac85x+1)$. (b) The normals are $e_\alpha$ and $R^{-1}e_\beta$; $e_\alpha=\pm R^{-1}e_\beta$ means $Re_\alpha=\pm e_\beta$, and for $\alpha\in\{\mu,\nu\}$ the vector $Re_\alpha$ has the components $\pm\frac35,\pm\frac45$ in the $(\mu,\nu)$-plane. (c) A bounded region with faces on hyperplanes needs face normals spanning $\R^3$; by (b) only $e_\rho$ is available. Both lattices descend to $\R^3/P$ exactly when $P\subset\Lk\cap R^{-1}\Lk$, which contains $50MJ\,\Z^3$ by \Cref{rot:lem:geom}(b).
\end{proof}

The ultraviolet analysis of \cite{K1,K2} needs boundaries that are lattice hyperplanes of the lattice being analysed (\cite[p.~326]{K2}: ``These hold only when the boundary is a lattice hyperplane orthogonal to a lattice coordinate direction. Clearly the two lattices in Fig.~1 never share such a boundary. Our solution of this problem is to use periodic boundary conditions.''). \Cref{rot:prop:forced} makes this precise: the objects on which the two lattices can be compared by King's method are tori with smooth sources, and smooth walls are the only admissible way to localise a torus; a Dirichlet region cannot be aligned with both lattices. The correlation inequalities give exactly ``cube $\le$ walled torus'' on the cone and never the reverse (\Cref{rot:lem:cutwall}; \Cref{rot:rem:gauss}), so the comparison ``walled torus $\le$ King's state $+\,\eta$'' must be supplied by an estimate on King's measure; in the curvature sector that estimate is \HS, through \Cref{rot:thm:shell}.

\begin{remark}[the content of \REa]\label{rot:rem:REScontent}
(i) At every fixed spacing \REa\ holds (\Cref{rot:prop:fixedeps}); its content is uniformity in the spacing. (ii) It is a decoupling statement for a heavy, weakly coupled region: on the shell the dimensionless couplings are of order $\lambda/\sqrt K$ and $e^2/\sqrt K$, and the plateau mass $n$ is free. The difficulty is not the shell but the comparison of two states that differ on the shell, seen through the interacting plateau and the light core, uniformly in the spacing. (iii) No correlation inequality bounds the tail from above: \Cref{lem:s02:BFS} gives signs and monotonicity, and the one comparison through a relative density that is uniform in $\eps$ is the Gaussian photon switch across a matter-free collar (\Cref{rot:lem:GB}); for the matter, the couplings across the boundary of the shell contribute to the logarithm of the relative density a term whose variance is of order $|\partial U|\eps^{-2}$. (iv) In the curvature sector the estimate is obtained from \HS\ in \Cref{rot:thm:shell}. The proof of \HS\ is a renormalisation-group expansion of the walled model of \cite{K1,K2} with the shell mass kept inside the matter covariance at every scale, run down to a stopping scale at which the plateau is massive in block units, the dependence on the shell mass being controlled through the $\varkappa$-Cauchy property of the reduced small-field action on a test set that does not depend on $\varkappa$. No such expansion for the symmetric massive phase of the three-dimensional model is in the printed literature; Ba{\l}aban, Imbrie and Jaffe \cite{BIJ88} develop a multiscale expansion of the high-frequency modes of the abelian Higgs model in the Higgs regime and postpone the large-field and decoupling estimates. \HS\ carries a window $n\ge n_*$ for the plateau height, which costs nothing for the theorem since $n$ is free. (v) At $e=0$ the photon is independent of the matter, and \REz\ holds trivially ($\tail_k\equiv0$ at $g=0$), while \REa\ for $g>0$ is a decoupling statement for the two-component $\varphi^4_3$ lattice model with a heavy shell. Nothing established at $e=0$ has been carried to $e>0$ in this problem, since King's mass counterterm at order $e^2$ is negative and of order $e^2/\eps$, so that every comparison of the $e>0$ model with an $e=0$ model at the bare level is a comparison with a model deep in its broken phase.
\end{remark}

\begin{remark}[the chain of the rotation part]\label{rot:rem:chain}
Every implication in the chain
\begin{gather*}
\HS\ \xrightarrow{\ \text{\Cref{rot:thm:shell}}\ }\ \REz\ \xrightarrow{\ \EC,\ \Kz,\ \DSh;\ \text{\Cref{rot:thm:REStoTIGHT}}\ }\ \TIz ,\\
\Kz+\TIz\ \xrightarrow{\ \text{\Cref{rot:thm:curv}}\ }\ \text{invariance under }\Rgen\ \xrightarrow{\ \text{\Cref{rot:thm:curv}}\ }\ O(3)\ \xrightarrow{\ \text{\Cref{cor:s06:TK}}\ }\ \Euc
\end{gather*}
is proved in this part, with \EC\ (\Cref{hyp:s03:EC}) wherever King's printed state is named (\Cref{rot:lem:free}, \Cref{rot:thm:curv}); \Cref{rot:thm:main} is the composition. Its two ends, the heavy-shell estimates \HS\ (\Cref{rot:hyp:HS}) and the torus comparison \Kz\ (\Cref{rot:hyp:K}), are renormalisation-group statements about King's walled torus model, proved in \Cref{rot:prop:SW1,rot:prop:SW3} and in \cref{ssec:rot:Kproof}; \DSh\ (\Cref{rot:hyp:DS}) is the dead-shell chain in the exterior-shell geometry.
\end{remark}
